\chapter{Moorings with MoorDyn}
\label{ch:moorings}

\section{Overview}
MoorDyn is a lumped-mass mooring line model~\cite{hall2015validation}. \fwt{} couples it to the
rigid-body motion as a \emph{restraint}: at every time step the body position and velocity are
passed to MoorDyn, which integrates the line dynamics and returns the mooring force and moment.
The HOS--OpenFOAM--MoorDyn coupling strategy and its validation are described
in~\cite{aliyar2022numerical}; the MoorDyn--OpenFOAM coupling follows Chen and
Hall~\cite{chen2022cfd}.

Two MoorDyn versions are built:
\begin{description}
  \item[MoorDyn v1] (\kw{moorDynR1}) lines between fixed anchors and fairleads on one body. Input
    in \file{Mooring/lines.txt}. MoorDyn v1 cannot store and restore its state, so the lines are
    advanced at most once per time step.
  \item[MoorDyn v2] (\kw{moorDynR2}, version 2.7.2) the current MoorDyn, with bodies, rods,
    bending stiffness, seabed friction and VTK output. Its state is saved and restored, so it
    can be called several times per time step (outer correctors, FloatStepper test steps).
\end{description}
Use MoorDyn v2 for new cases. The MoorDyn input file formats are described in the MoorDyn
documentation (\url{https://moordyn.readthedocs.io}).

Both versions are available for both body-motion solvers:
\begin{center}\small
\begin{tabular}{@{}lll@{}}
\toprule
Solvers & Restraint keyword & Library in \kw{libs} \\
\midrule
\app{floatStepperWaveIsoFoam}, & \kw{floaterMotionRestraint} & \file{moorDynForFloatStepper} \\
\quad\app{floatStepper} & & \\[2pt]
\app{waveFoam}, \app{waveIsoFoam} & \kw{sixDoFRigidBodyMotionRestraint} & \file{moorDynR1SixDoF} (v1) \\
\quad(\kw{sixDoFRigidBodyMotion}) & & \file{moorDynR2SixDoF} (v2) \\
\bottomrule
\end{tabular}
\end{center}

\section{Moorings with FloatStepper}
Add the library in \file{system/controlDict} and the restraint in the \kw{restraints}
sub-dictionary of Code Frame~\ref{lst:fsdyn}:
\begin{lstlisting}[caption={MoorDyn v2 restraint for FloatStepper.},label={lst:fsmoor}]
// system/controlDict
libs    (moorDynForFloatStepper);

// constant/dynamicMeshDict, inside floaterMeshMotionSolverCoeffs
restraints
{
    moorDynR2_body
    {
        floaterMotionRestraint  moorDynR2;
        couplingMode    "BODY";
        inputFile       "Mooring/lines_v2_bodyC.txt";
        writeMooringVTK true;
        vtkPrefix       "mdv2_body";
        vtkStartTime    0;
        outerCorrector  1;
    }
}
\end{lstlisting}
For MoorDyn v1, the restraint has no further entries; the lines are read from
\file{Mooring/lines.txt}:
\begin{lstlisting}
moorDynR1
{
    floaterMotionRestraint  moorDynR1;
}
\end{lstlisting}

\section{Moorings with \texorpdfstring{\kw{sixDoFRigidBodyMotion}}{sixDoFRigidBodyMotion}}
\label{sec:moordyn-sixdof}
For \app{waveFoam} and \app{waveIsoFoam}, the same restraints are available for OpenFOAM's
\kw{sixDoFRigidBodyMotion}, so a moored body can be simulated without FloatStepper:
\begin{lstlisting}[caption={MoorDyn v2 restraint for \kw{sixDoFRigidBodyMotion}.},label={lst:sixdofmoor}]
// system/controlDict
libs    ( "libwaves2Foam.so" moorDynR2SixDoF );

// constant/dynamicMeshDict
dynamicFvMesh       dynamicMotionSolverFvMesh;
motionSolverLibs    (sixDoFRigidBodyMotion);
motionSolver        sixDoFRigidBodyMotion;

sixDoFRigidBodyMotionCoeffs
{
    patches         (box);
    innerDistance   0.075;
    outerDistance   0.4;
    mass            3.148;
    centreOfMass    (0 0 -0.0126);
    momentOfInertia (0.015 0.0167 0.021);
    accelerationRelaxation 0.4;
    solver          { type Newmark; }
    constraints     { ... }

    restraints
    {
        moorDynR2_body
        {
            sixDoFRigidBodyMotionRestraint  moorDynR2;
            couplingMode    "BODY";
            inputFile       "Mooring/lines_v2_bodyC.txt";
            writeMooringVTK true;
            vtkPrefix       "mdv2_body";
        }
    }
}
\end{lstlisting}
For MoorDyn v1 use \kw{sixDoFRigidBodyMotionRestraint moorDynR1;} with the library
\file{moorDynR1SixDoF}. When the solver uses several outer correctors with
\kw{moveMeshOuterCorrectors}, \kw{moorDynR1} returns the force of the first call in each time
step; \kw{moorDynR2} restores the MoorDyn state and recomputes it.

\section{Restraint options}
\begin{table}[htbp]
\centering\small
\begin{tabularx}{\linewidth}{@{}l l X@{}}
\toprule
Keyword & Default & Meaning \\
\midrule
\kw{inputFile}        & -- (required) & MoorDyn v2 input file, relative to the case folder \\
\kw{couplingMode}     & -- (required) & \kw{"BODY"}: MoorDyn has one coupled body (6 degrees of freedom);
                                        \kw{"POINT"}: MoorDyn has coupled points \\
\kw{refAttachmentPt}  & -- & list of fairlead points, one per coupled point (\kw{POINT} only) \\
\kw{writeMooringVTK}  & \kw{false} & write the line geometry as VTK files for ParaView \\
\kw{vtkPrefix}        & \kw{mdV2} & file name prefix of the VTK files \\
\kw{vtkStartTime}     & 0 & first time at which VTK files are written \\
\kw{vtkLegacyFormat}  & (build-dependent) & legacy \file{.vtk} instead of XML VTK \\
\kw{outerCorrector}   & 1 & outer-corrector iteration in which the VTK files are written \\
\kw{allowStaticRestart} & \kw{false} & on a restart without a saved MoorDyn state, re-initialise the
                                        lines from static equilibrium (warning) instead of stopping;
                                        also accepted by \kw{moorDynR1} \\
\bottomrule
\end{tabularx}
\caption{Options of the \kw{moorDynR2} restraint (FloatStepper and \kw{sixDoFRigidBodyMotion}).}
\end{table}

In \kw{BODY} mode, the MoorDyn input file must define exactly one coupled body; its position is
taken from the rigid-body motion. In \kw{POINT} mode, the number of entries in
\kw{refAttachmentPt} must equal the number of coupled points in the MoorDyn input file.

\section{Restart}
At every write time, \kw{moorDynR2} saves the complete MoorDyn state (line nodes, coupled body,
time integrator) with the time folder, in
\file{<time>/uniform/moorDyn/<restraint name>.dat} (\file{processor0/<time>/...} in parallel;
\file{reconstructPar} copies it to the case folder and \file{decomposePar} links to it). When the
solver restarts from a written time, the lines continue from this state instead of a new static
equilibrium, so the mooring force of the first step equals that of the continuous run. This holds
for \file{waveIsoFoam}, \file{waveIsoFoamSWENSE} and \file{floatStepperWaveIsoFoam}. A restart from
a time folder without this file (for example, one written before this feature) stops with an error,
unless \kw{allowStaticRestart true} is set. MoorDyn v1 cannot save its state, so \kw{moorDynR1}
always stops on a restart unless \kw{allowStaticRestart true} is set; its lines then restart from
static equilibrium.

MoorDyn rewrites its \file{.out} files when it starts. On a restart, \kw{moorDynR2} therefore
first keeps their rows up to the restart time in \file{<file>.out.prev}; at the end of the run
they are merged back in front of the new rows and the \file{.prev} file is removed, so each
\file{.out} file covers the whole simulation. If a restarted run stops early (crash, wall-time
limit), the \file{.prev} files stay and the next restart merges them. Rows of the old run after the
restart time are dropped. VTK file numbering continues after the last frame before the restart.
The MoorDyn log file (\file{.log}) is started again.

\section{Output}
MoorDyn writes the line tensions and node positions to the \file{Mooring/} folder (MoorDyn v1:
\file{Lines.out} and \file{Line<n>.out}; MoorDyn v2: as set in its input file). With
\kw{writeMooringVTK true}, the lines are also written as VTK files that can be loaded in
ParaView together with the CFD solution; \file{tutorials/floatStepperWaveIsoFoam} contains the script
\file{MooringVisualizationInParaview.py} for this.
